\documentclass{article}
\usepackage[utf8]{inputenc}
\usepackage[ngerman]{babel}
\usepackage{amsmath, amssymb, amsthm}
\usepackage{geometry}

\usepackage{enumitem}
\setlist[enumerate]{leftmargin=1cm, itemsep=2pt}

\geometry{a4paper, margin=2.5cm}
\newtheorem{satz}{Satz}[section]
\newtheorem{axiom}[satz]{Axiom}
\newtheorem{lemma}[satz]{Lemma}
\newtheorem{theorem}[satz]{Theorem}
\theoremstyle{definition}
\newtheorem{definition}[satz]{Definition}
\newtheorem{bemerkung}[satz]{Bemerkung}
\begin{document}
\title{Wiki Export}
\author{LLM-Wiki Parser}
\maketitle

% --- darboux_integrale_definition ---
\begin{definition}[Darbouxsche Ober- und Unterintegrale]
Für eine beschränkte Funktion $f$ auf $[a, b]$ und eine Zerlegung $Z$ definiert man die \textbf{Untersumme} $U(f, Z)$ und die \textbf{Obersumme} $O(f, Z)$ durch die Infima bzw. Suprema der Funktionswerte in den Teilintervallen.
Das \textbf{Darbouxsche Unterintegral} ist das Supremum aller Untersummen:
$$\underline{\int_a^b} f(x) \, dx = \sup_{Z} U(f, Z)$$
Das \textbf{Darbouxsche Oberintegral} ist das Infimum aller Obersummen:
$$\overline{\int_a^b} f(x) \, dx = \inf_{Z} O(f, Z)$$
Eine Funktion heißt \textbf{D-integrierbar}, wenn Ober- und Unterintegral übereinstimmen.


\end{definition}
\begin{proof}
Da für jede Zerlegung $U(f, Z) \le O(f, Z)$ gilt, ist stets das Unterintegral kleiner oder gleich dem Oberintegral. Die Übereinstimmung beider Werte ist äquivalent zur Riemann-Integrierbarkeit.
\end{proof}
% --- stammfunktion_definition ---
\begin{definition}[Definition der Stammfunktion]
Eine Funktion $F$ heißt \textbf{Stammfunktion} von $f$ auf einem Intervall $I$, wenn $F$ auf $I$ differenzierbar ist und für alle $x \in I$ gilt:
$$F'(x) = f(x)$$


\end{definition}
\begin{proof}
Heuser (Nr. 55): Ist $F_0$ eine Stammfunktion von $f$, so sind alle Funktionen der Form $F(x) = F_0(x) + C$ mit einer Konstanten $C \in \mathbb{R}$ ebenfalls Stammfunktionen von $f$. Dies folgt daraus, dass die Ableitung einer Konstanten null ist und Funktionen mit verschwindender Ableitung auf einem Intervall konstant sind (Folgerung aus dem Mittelwertsatz).
\end{proof}
% --- riemann_integral_definition ---
\begin{definition}[Definition des Riemann-Integrals]
Eine auf dem kompakten Intervall $[a, b]$ erklärte reelle Funktion $f$ heißt \textbf{Riemann-integrierbar} (kurz: R-integrierbar), wenn für jede Folge von Zerlegungen $(Z_n)$ mit einer Feinheit $|Z_n| \to 0$ die zugehörigen \textbf{Riemannschen Summen}
$$S(Z_n, \xi) = \sum_{k=1}^n f(\xi_k) \Delta x_k$$
(unabhängig von der Wahl der Zwischenstellen $\xi_k$) gegen denselben Grenzwert $J$ konvergieren. Dieser Grenzwert $J$ heißt das \textbf{Riemann-Integral} von $f$ über $[a, b]$ und wird bezeichnet als:
$$\int_a^b f(x) \, dx$$


\end{definition}
\begin{proof}
Die Definition stellt sicher, dass das Integral als Grenzwert eines Netzes von Summen existiert. Eine notwendige Bedingung für die R-Integrierbarkeit ist die Beschränktheit der Funktion $f$ auf $[a, b]$.
\end{proof}
% --- uneigentliches_integral_unbeschraenkt_definition ---
\begin{definition}[Uneigentliche Integrale über unbeschränkte Intervalle]
Ein Integral über ein unbeschränktes Intervall wird als Grenzwert von Riemann-Integralen über endlichen Teilintervallen definiert.

1. Sei $f$ für jedes $t > a$ auf $[a, t]$ R-integrierbar. Existiert der Grenzwert
   $$\int_{a}^{\infty} f(x) \, dx := \lim_{t \to \infty} \int_{a}^{t} f(x) \, dx$$
   als endliche Zahl, so heißt das uneigentliche Integral \textbf{konvergent}. Andernfalls heißt es \textbf{divergent}.

2. Entsprechend definiert man für eine auf $[t, b]$ integrierbare Funktion:
   $$\int_{-\infty}^{b} f(x) \, dx := \lim_{t \to -\infty} \int_{t}^{b} f(x) \, dx$$

3. Ein Integral über $(-\infty, \infty)$ heißt konvergent, wenn für ein beliebiges $a \in \mathbb{R}$ beide Integrale $\int_{-\infty}^{a} f(x) \, dx$ und $\int_{a}^{\infty} f(x) \, dx$ konvergieren. Man setzt dann:
   $$\int_{-\infty}^{\infty} f(x) \, dx := \int_{-\infty}^{a} f(x) \, dx + \int_{a}^{\infty} f(x) \, dx$$


\end{definition}
\begin{proof}
Die Definition erweitert den Integralbegriff durch einen Grenzübergang. Die Wahl von $a$ im dritten Fall ist unerheblich, da für $a < b$ gilt:
$$\int_{t}^{a} f(x) \, dx + \int_{a}^{\infty} f(x) \, dx = \int_{t}^{b} f(x) \, dx - \int_{a}^{b} f(x) \, dx + \int_{a}^{\infty} f(x) \, dx$$
Beim Grenzübergang $t \to -\infty$ heben sich die konstanten Anteile konsistent auf.
\end{proof}
% --- integrationsregeln_arithmetik ---
\begin{satz}[Arithmetische Regeln für Integrale]
Sind $f$ und $g$ auf $[a, b]$ R-integrierbar, so gilt:
\begin{enumerate}[label=\arabic*.]
    \item \textbf{Linearität}: $\int_a^b (\alpha f + \beta g) \, dx = \alpha \int_a^b f \, dx + \beta \int_a^b g \, dx$ für $\alpha, \beta \in \mathbb{R}$.
    \item \textbf{Monotonie}: Ist $f(x) \le g(x)$ für alle $x \in [a, b]$, so ist $\int_a^b f \, dx \le \int_a^b g \, dx$.
    \item \textbf{Intervalladditivität}: $\int_a^b f \, dx = \int_a^c f \, dx + \int_c^b f \, dx$ für $c \in (a, b)$.
\end{enumerate}


\end{satz}
\begin{proof}
Diese Eigenschaften folgen unmittelbar aus den entsprechenden Eigenschaften der endlichen Riemannschen Summen durch Grenzübergang.
\end{proof}
% --- mittelwertsatz_integralrechnung ---
\begin{satz}[Erster Mittelwertsatz der Integralrechnung]
Ist $f$ auf $[a, b]$ stetig, so gibt es mindestens eine Stelle $\xi \in [a, b]$, so dass gilt:
$$\int_a^b f(x) \, dx = f(\xi) \cdot (b - a)$$


\end{satz}
\begin{proof}
Da $f$ auf dem kompakten Intervall $[a, b]$ stetig ist, nimmt sie nach dem Satz von Weierstraß ein Minimum $m$ und ein Maximum $M$ an. Es gilt $m(b-a) \le \int f \, dx \le M(b-a)$. Aus dem Zwischenwertsatz folgt dann die Existenz von $\xi$.
\end{proof}
% --- erweiterter_mittelwertsatz_integralrechnung ---
\begin{satz}[Erweiterter Mittelwertsatz der Integralrechnung]
Es seien $f, g$ auf $[a, b]$ R-integrierbar, $f$ stetig und $g$ ohne Vorzeichenwechsel (d.h. $g(x) \ge 0$ oder $g(x) \le 0$ für alle $x$). Dann gibt es eine Stelle $\xi \in [a, b]$ mit:
$$\int_a^b f(x)g(x) \, dx = f(\xi) \int_a^b g(x) \, dx$$


\end{satz}
\begin{proof}
Der Beweis verläuft analog zum ersten Mittelwertsatz unter Ausnutzung der Monotonie des Integrals, wobei die Funktion $g$ als "Gewichtsfunktion" fungiert.
\end{proof}
% --- hauptsatz_differential_integralrechnung ---
\begin{satz}[Hauptsatz der Differential- und Integralrechnung]
Der Hauptsatz besteht aus zwei komplementären Aussagen:

\begin{enumerate}[label=\arabic*.]
    \item \textbf{Teil (Existenz von Stammfunktionen)}: Ist $f: [a, b] \to \mathbb{R}$ stetig, so ist die durch
    $$F(x) = \int_a^x f(t) \, dt$$
    definierte Integralfunktion eine Stammfunktion von $f$ auf $[a, b]$, d.h. $F'(x) = f(x)$.

    \item \textbf{Teil (Berechnung von Integralen)}: Ist $F$ eine beliebige Stammfunktion der R-integrierbaren Funktion $f$ auf $[a, b]$, so gilt:
    $$\int_a^b f(x) \, dx = F(b) - F(a) = [F(x)]_a^b$$
\end{enumerate}


\end{satz}
\begin{proof}
Der Beweis des ersten Teils nutzt den Mittelwertsatz der Integralrechnung, um zu zeigen, dass der Differenzenquotient von $F$ gegen $f(x)$ strebt. Der zweite Teil folgt daraus, dass sich zwei Stammfunktionen nur um eine Konstante unterscheiden.
\end{proof}
% --- integralkriterium_konvergenz ---
\begin{satz}[Integralkriterium für unendliche Reihen]
Das Integralkriterium stellt einen Zusammenhang zwischen der Konvergenz einer unendlichen Reihe und einem uneigentlichen Integral her.

Die Funktion $f$ sei auf $[m, \infty)$ (mit $m \in \mathbb{N}$) positiv und monoton fallend. Dann haben die Reihe $\sum_{k=m}^{\infty} f(k)$ und das uneigentliche Integral $\int_{m}^{\infty} f(x) \, dx$ dasselbe Konvergenzverhalten.

Es gilt insbesondere:
$$\sum_{k=m+1}^{\infty} f(k) \le \int_{m}^{\infty} f(x) \, dx \le \sum_{k=m}^{\infty} f(k)$$


\end{satz}
\begin{proof}
Da $f$ fallend ist, gilt für $x \in [k, k+1]$ die Abschätzung:
$$f(k+1) \le f(x) \le f(k)$$
Integration über $[k, k+1]$ liefert:
$$f(k+1) \le \int_{k}^{k+1} f(x) \, dx \le f(k)$$
Summation über $k = m, \dots, n-1$ ergibt:
$$\sum_{k=m+1}^{n} f(k) \le \int_{m}^{n} f(x) \, dx \le \sum_{k=m}^{n-1} f(k)$$
Durch den Grenzübergang $n \to \infty$ folgt die Behauptung, da beide Seiten monoton wachsen und somit genau dann konvergieren, wenn sie beschränkt sind.
\end{proof}
% --- stetigkeit_impliziert_integrabilitaet ---
\begin{satz}[Integrierbarkeit stetiger Funktionen]
Jede auf einem kompakten Intervall $[a, b]$ \textbf{stetige} Funktion $f$ ist dort Riemann-integrierbar.


\end{satz}
\begin{proof}
Da $f$ auf dem kompakten Intervall $[a, b]$ stetig ist, ist sie dort auch gleichmäßig stetig (Satz von Heine). Dies erlaubt es, die Differenz zwischen Ober- und Untersummen (Darboux-Integrale) durch eine beliebig kleine Zahl $\epsilon$ abzuschätzen, sobald die Feinheit der Zerlegung klein genug ist.
\end{proof}
% --- uneigentliches_integral_konvergenzkriterien ---
\begin{satz}[Konvergenzkriterien für uneigentliche Integrale]
Für die Konvergenz uneigentlicher Integrale $\int_{a}^{\infty} f(x) \, dx$ gelten folgende Kriterien:

1. \textbf{Cauchy-Kriterium}: Das Integral konvergiert genau dann, wenn zu jedem $\epsilon > 0$ ein $s_0$ existiert, so dass für alle $t > s > s_0$ gilt:
   $$\left| \int_{s}^{t} f(x) \, dx \right| < \epsilon$$

2. \textbf{Monotoniekriterium}: Ist $f \ge 0$ auf $[a, \infty)$, so konvergiert $\int_{a}^{\infty} f(x) \, dx$ genau dann, wenn die Teilintegrale $\int_{a}^{t} f(x) \, dx$ nach oben beschränkt sind.

3. \textbf{Majorantenkriterium}: Gilt $|f| \le g$ auf $[a, \infty)$ und konvergiert $\int_{a}^{\infty} g(x) \, dx$, so konvergiert $\int_{a}^{\infty} f(x) \, dx$ (sogar absolut).

4. \textbf{Minorantenkriterium}: Gilt $0 \le h \le f$ und divergiert $\int_{a}^{\infty} h(x) \, dx$, so divergiert auch $\int_{a}^{\infty} f(x) \, dx$.

5. \textbf{Grenzwertkriterium}: Sind $f, g > 0$ und existiert der Grenzwert $\lim_{x \to \infty} \frac{f(x)}{g(x)} = L$ mit $L \in (0, \infty)$, so haben beide Integrale dasselbe Konvergenzverhalten.


\end{satz}
\begin{proof}
Diese Kriterien sind direkte Übertragungen der Kriterien für Zahlenfolgen und unendliche Reihen auf die Funktion $F(t) = \int_{a}^{t} f(x) \, dx$. Die absolute Konvergenz (Punkt 3) folgt aus der Vollständigkeit von $\mathbb{R}$.
\end{proof}
% --- partielle_integration ---
\begin{satz}[Partielle Integration]
Sind $u$ und $v$ auf $[a, b]$ stetig differenzierbare Funktionen, so gilt:
$$\int_a^b u(x) v'(x) \, dx = [u(x)v(x)]_a^b - \int_a^b u'(x) v(x) \, dx$$


\end{satz}
\begin{proof}
Die Regel folgt durch Integration der Produktregel $(uv)' = u'v + uv'$ und Anwendung des Hauptsatzes der Differential- und Integralrechnung.
\end{proof}
% --- riemann_integrabilitaet_kriterium ---
\begin{satz}[Riemannsches Integrabilitätskriterium]
Eine beschränkte Funktion $f$ ist auf $[a, b]$ genau dann Riemann-integrierbar, wenn zu jedem $\epsilon > 0$ eine Zerlegung $Z$ von $[a, b]$ existiert, so dass die Differenz zwischen Obersumme und Untersumme kleiner als $\epsilon$ ist:
$$O(f, Z) - U(f, Z) < \epsilon$$


\end{satz}
\begin{proof}
Das Kriterium folgt direkt aus der Definition der Darboux-Integrale als Supremum bzw. Infimum. Es ist das praktische Hauptwerkzeug, um die Integrierbarkeit von Funktionenklassen (wie stetigen oder monotonen Funktionen) nachzuweisen.
\end{proof}
\end{document}
